% Part: lambda-calculus
% Chapter: lambda-definability
% Section: lambda-definable-recursive

\documentclass[../../../include/open-logic-section]{subfiles}

\begin{document}

\olfileid{lam}{dfl}{ldr}
\olsection{\usetoken{S}{lambda definable} تابعې بازګشتي دي}

نۀ يوازې ټولې جزوي بازګشتي تابعې !!{lambda definable}
دي، بلکې برعکس خبره هم رښتيا ده: ټولې
!!{lambda definable} تابعې جزوي بازګشتي دي۔

\begin{thm}
  \ollabel{thm:lambda-computable} که يوه جزوي تابع~$f$
  !!{lambda definable} وي، نو جزوي بازګشتي هم ده۔
\end{thm}

\begin{proof}
  د ثبوت يوازې لنډه طرحه وړاندې کوو۔ لومړے د
  $\lambd$-ترمونو حسابي کول کوو: د لړيو د کوډولو
  دپاره د اوليو عددونو د توانونو په معمول طريقه هر
  $\lambd$-ترم ته په منظم ډول د ګوډل شمېره ورکوو۔
  بيا يوه جزوي بازګشتي تابع $\fn{normalize}(t)$
  تعريفوو چې د لامبډا ترم د ګوډل شمېره~$t$ د
  آرګومېنټ په توګه اخلي؛ که ترم عادي بڼه ولري، د
  هغې بڼې د ګوډل شمېره ورکوي، او که نۀ وي،
  ناتعريفه پاتې کېږي۔ ورپسې دوه جزوي بازګشتي
  تابعې $\fn{toChurch}$ او $\fn{fromChurch}$
  تعريفوو چې طبيعي عددونه د اړوندو چرچ عددنښو
  د ګوډل شمېرو ته اړوي او له هغو شمېرو ئې بېرته
  طبيعي عددونو ته رااړوي۔

  له دغو بازګشتي تابعو څخه په کار اخيستو، تابع~$f$
  د جزوي بازګشتي تابعې په توګه تعريفولے شو۔ داسې
  $\lambd$-ترم~$F$ شته چې~$f$ !!{lambda define}s۔
  د $f(n_1, \dots, n_k)$ د محاسبې دپاره لومړے د
  $\fn{toChurch}(n_i)$ په وسيله د اړوندو چرچ
  عددنښو د ګوډل شمېرې ترلاسه کوو۔ بيا دا شمېرې
  له $\Gn{F}$ سره يوځاے کوو، څو د
  $F \num{n_1}\dots\num{n_k}$ ترم د ګوډل شمېره
  ترلاسه شي۔ اوس پر دې ګوډل شمېره
  $\fn{normalize}$ کاروو۔ که $f(n_1, \dots,
  n_k)$ تعريف شوې وي، $F \num{n_1}\dots\num{n_k}$
  عادي بڼه لري، او هغه بايد د چرچ عددنښه وي۔ که
  تابع ناتعريفه وي، ترم عادي بڼه
  نۀ لري؛ نو
  \[\fn{normalize}(\Gn{F\num{n_1}\dots\num{n_k}})\]
  هم ناتعريفه وي۔ په پاې کښې د عادي شوې بڼې د
  ګوډل پر شمېره $\fn{fromChurch}$ کاروو۔
\end{proof}

\end{document}
